% Einheit 3 von 3 – Strahlensätze (bank/strahlensaetze/e3.jsonl, zusammenbau v0.1)
%% TODO Zweigzeile Teil 1: was hier gelernt wird (Satz in Schülersprache aus dem Einheitstitel) – Einheit 3
\einheitenkopf[e3]{Einheit 3 von 3 · Strahlensätze}
\zweigzeile{ab Kl. 9 · keine P10-Aufgabe}
\setcounter{aufgabe}{18}

%% TODO Titel: Ich-kann-Satz fehlt in der Bank (Platzhalter: Kettenname) – Nr. 19
%% TODO Anweisung: ein Satz über der ersten Teilaufgabe fehlt in der Bank – Nr. 19
\begin{aufgabe}{Strahlensätze}
\begin{teile}
\teil Die Grafik zeigt zwei Geraden durch $Z$ und zwei Parallelen. V-Figur oder X-Figur? \kreuz{V-Figur} \\ \kreuz{X-Figur}
\end{teile}
\swz{In der V-Figur ist $\overline{ZA} = 3\,\text{cm}$, $\overline{ZA'} = 6\,\text{cm}$ und $\overline{ZB} = 4\,\text{cm}$. Berechne $\overline{ZB'}$.}{\strahlensatz{3}{6}{35}}
\swz{In der V-Figur ist $\overline{ZA} = 2\,\text{cm}$, $\overline{ZA'} = 8\,\text{cm}$ und $\overline{ZB} = 5\,\text{cm}$. Berechne $\overline{ZB'}$.}{\strahlensatz{2}{7}{35}}
\swz{In der V-Figur ist $\overline{ZA} = 5\,\text{cm}$, $\overline{ZA'} = 15\,\text{cm}$ und $\overline{ZB} = 2\,\text{cm}$. Berechne $\overline{ZB'}$.}{\strahlensatz{3}{7}{35}}
\swz{In der V-Figur ist $\overline{ZA} = 3\,\text{cm}$, $\overline{ZA'} = 9\,\text{cm}$ und $\overline{ZB} = 5\,\text{cm}$. Berechne $\overline{ZB'}$.}{\strahlensatz{3}{7}{35}}
\swz{In der V-Figur ist $\overline{ZA} = 4\,\text{cm}$, $\overline{ZA'} = 16\,\text{cm}$ und $\overline{ZB} = 3\,\text{cm}$. Berechne $\overline{ZB'}$.}{\strahlensatz{3}{7}{35}}
\swfrage{Was bleibt gleich, was ändert sich?}
\swz{In der V-Figur ist $\overline{ZA} = 2\,\text{cm}$, $\overline{ZA'} = 7\,\text{cm}$ und $\overline{AB} = 4\,\text{cm}$. Berechne die Parallele $\overline{A'B'}$.}{\strahlensatz{2}{7}{35}}
\swz{In der X-Figur liegen $A$ und $A'$ auf verschiedenen Seiten von $Z$. Es ist $\overline{ZA} = 2\,\text{cm}$, $\overline{ZA'} = 3\,\text{cm}$ und $\overline{ZB} = 4\,\text{cm}$. Berechne $\overline{ZB'}$.}{\strahlensatz[x]{2}{3}{30}}
\swz{In der V-Figur ist $\overline{ZA} = 2{,}4\,\text{cm}$, $\overline{ZA'} = 6\,\text{cm}$ und $\overline{ZB} = 3{,}2\,\text{cm}$. Stelle eine Verhältnisgleichung auf und berechne $\overline{ZB'}$.}{\strahlensatz{2.4}{6}{35}}
\swz{In der V-Figur ist $\overline{ZA} = 4\,\text{cm}$, $\overline{AA'} = 3\,\text{cm}$ und $\overline{ZB} = 6\,\text{cm}$. Berechne $\overline{BB'}$.}{\strahlensatz{2}{3.5}{35}}
\swz{Ein Stab ist $1{,}5\,\text{m}$ hoch und wirft einen $2\,\text{m}$ langen Schatten. Er steht so, dass sein Schatten genau dort endet, wo der Schatten eines Baums endet; dieser ist $14\,\text{m}$ lang. In der Skizze (nicht maßstabsgerecht) ist $Z$ das Schattenende, $A$ und $B$ Fuß und Spitze des Stabs, $A'$ und $B'$ Fuß und Spitze des Baums. Wie hoch ist der Baum?}{\strahlensatz{1.5}{6}{30}}
\swz[3]{Zwei Bäume am anderen Ufer eines Flusses stehen $24\,\text{m}$ auseinander. Von einem Punkt $Z$ aus werden sie angepeilt; die zwei Peillinien treffen das eigene Ufer $8\,\text{m}$ auseinander. $Z$ ist $5\,\text{m}$ vom eigenen Ufer entfernt, beide Ufer sind parallel. Fertige eine Skizze an und berechne die Breite des Flusses.}{\rechenplatz[halb]{4}}
%% TODO Darstellung für schwach fehlt in der Bank (strahlensaetze-e3-k1-s8-v1; Abschnitt „Für schwache Schüler“)
\swz[3]{In einer V-Figur ist $\overline{ZA} = 4\,\text{cm}$, $\overline{ZA'} = 7\,\text{cm}$, $\overline{ZB} = 8\,\text{cm}$ und $\overline{ZB'} = 14\,\text{cm}$. Sind die Geraden $AB$ und $A'B'$ parallel? Vergleiche die Verhältnisse.}{}
\swz[3]{Ein $1{,}80\,\text{m}$ großer Mann wirft einen $2{,}40\,\text{m}$ langen Schatten. Zur selben Zeit ist der Schatten eines Hochhauses $64\,\text{m}$ lang. Fertige eine Skizze an, stelle eine Verhältnisgleichung auf und berechne die Höhe des Hochhauses. Schreibe einen Antwortsatz.}{\rechenplatz[halb]{4}}
\end{aufgabe}

%% TODO Titel: Ich-kann-Satz fehlt in der Bank (Platzhalter: Kettenname) – Nr. 20
%% TODO Anweisung: ein Satz über der ersten Teilaufgabe fehlt in der Bank – Nr. 20
\begin{aufgabe}{Voraussetzung prüfen}
\begin{teile}
\teil In einer V-Figur bilden $AB$ und $A'B'$ mit dem Strahl durch $Z$, $A$ und $A'$ beide einen Winkel von $65^\circ$ (Stufenwinkel). Darf man den Strahlensatz anwenden? \kreuz{ja} \\ \kreuz{nein}
\end{teile}
\end{aufgabe}

%% TODO Titel: Ich-kann-Satz fehlt in der Bank (Platzhalter: Kettenname) – Nr. 21
%% TODO Anweisung: ein Satz über der ersten Teilaufgabe fehlt in der Bank – Nr. 21
\begin{aufgabe}{Strecke in n gleiche Teile teilen}
\swa{Zeichne eine Strecke von $7\,\text{cm}$ Länge und teile sie mit einem Hilfsstrahl und Parallelen in $3$ gleiche Teile. Wie lang ist ein Teil?}{\begin{ksys}[xmin=0,xmax=12,ymin=0,ymax=6]\end{ksys}}
\end{aufgabe}

%% TODO Titel: Ich-kann-Satz fehlt in der Bank (Platzhalter: Kettenname) – Nr. 22
%% TODO Anweisung: ein Satz über der ersten Teilaufgabe fehlt in der Bank – Nr. 22
\begin{aufgabe}{Strahlensätze – Fehler finden, Begründen, Anwendung, Darstellungswechsel}
\swz[3]{In einer V-Figur ist $\overline{ZA} = 4\,\text{cm}$, $\overline{AA'} = 2\,\text{cm}$ und $\overline{AB} = 3\,\text{cm}$. Jan rechnet: $\overline{A'B'} : 3 = 2 : 4$, also $\overline{A'B'} = 1{,}5\,\text{cm}$. Finde den Fehler und rechne richtig.}{}
\swfrage{Tim sagt: „Der Strahlensatz gilt auch, wenn die zwei Linien nur fast parallel sind.“ Stimmt das? Begründe.}
\swz[3]{Eine Leiter lehnt an einer Wand und reicht $4\,\text{m}$ hoch; ihr Fuß steht $1{,}6\,\text{m}$ vor der Wand. Ein Maler steht auf der Leiter in $2{,}5\,\text{m}$ Höhe. Kann er mit ausgestrecktem Arm ($70\,\text{cm}$) die Wand erreichen?}{}
\begin{teile}
\teil In der V-Figur sind $AB$ und $A'B'$ parallel. Welche Gleichung passt? \kreuz{$\overline{ZA} : \overline{ZA'} = \overline{ZB} : \overline{ZB'}$} \\ \kreuz{$\overline{ZA} : \overline{AA'} = \overline{ZB'} : \overline{ZB}$} \\ \kreuz{$\overline{AB} : \overline{A'B'} = \overline{ZA} : \overline{AA'}$}
\end{teile}
\end{aufgabe}

\uebersichtskasten{Strahlensatzfigur: Zwei Geraden durch ein Zentrum Z und zwei Parallelen – beim V liegen beide Parallelen auf derselben Seite von Z, beim X auf verschiedenen Seiten. Die Sätze gelten nur, wenn die zwei Linien wirklich parallel sind. \\ Erster Strahlensatz: Die Abschnitte auf den beiden Strahlen, vom Zentrum aus gemessen, stehen im selben Verhältnis: ZA : ZA′ = ZB : ZB′. ZA = 2 cm, ZA′ = 6 cm, ZB = 3 cm → ZB′ = 9 cm (dreimal so lang). \\ Zweiter Strahlensatz: Die Parallelen verhalten sich wie die Abschnitte vom Zentrum bis zu ihnen: AB : A′B′ = ZA : ZA′. AB = 4 cm → A′B′ = 12 cm. \\ Rechnen: Verhältnisgleichung aufstellen, die gesuchte Strecke allein auf eine Seite bringen, dann ausrechnen: x : 6 = 3 : 2 → x = 6 · 3 : 2 = 9.}
